% Cover Letter: Telnyx Senior Software Engineer, Java
% Author: Ajay Venkatesh

\documentclass[letterpaper,11pt]{article}

\usepackage[top=0.8in, bottom=0.8in, left=0.9in, right=0.9in]{geometry}
\usepackage[hidelinks]{hyperref}
\usepackage{lmodern}
\usepackage{parskip}
\usepackage{microtype}

\setlength{\parskip}{6pt}
\setlength{\parindent}{0pt}

\begin{document}

\begin{flushleft}
\textbf{\large Ajay Venkatesh} \\
Brooklyn, NY \\
+1 516 585 9013 \\
\href{mailto:av3855@nyu.edu}{av3855@nyu.edu}
\end{flushleft}

\vspace{10pt}

May 20, 2026

\vspace{6pt}

Hiring Manager \\
Telnyx

\vspace{10pt}

Dear Hiring Manager,

My name is Ajay Venkatesh. I'm finishing my M.S. in Computer Engineering at NYU, with several years of production software experience: a few years at PwC in Bengaluru, a summer SDE internship at Amazon, and ongoing on-campus work at NYU's Novel AI Technologies. My focus has been on \textbf{Java backend services}, distributed message-driven systems, and shipping infrastructure that has to stay up.

The Telnyx posting reads like a description of the kind of work I want to do next. At Amazon I engineered backend services in \textbf{Java} (Coral, Smithy, Dagger, CBOR) with low-level design patterns and unit/integration testing via \textbf{Mockito}, shipped under code-review discipline and frequent CI/CD deployments. On top of that I built a high-throughput, event-driven processing system on \textbf{AWS ECS Fargate, SQS, SNS}, decoupling distributed message handling and materially reducing response latency under variable load. I provisioned resilient \textbf{Infrastructure-as-Code} on \textbf{AWS CDK} across auto-scaling clusters with multi-stage, multi-region CI/CD and rollback controls. The shape of that work, scaling messaging-style services to stay highly available across regions, is the same shape your 30+ microservices need to keep running.

On the data-pipeline and streaming side, I'm fluent in \textbf{Kafka} and event-driven architectures, and I've shipped \textbf{Kubernetes} and \textbf{Docker}-based deployments in production. At PwC I spent three years on a distributed backend on \textbf{Microsoft Azure} with event-driven processing, fault-tolerant retries over \textbf{SQL Server}, and parallelized batch ingestion that cut execution from hours to minutes. I'm comfortable across \textbf{PostgreSQL, MySQL, MongoDB, DynamoDB,} and \textbf{Cassandra} from coursework, and I learn new datastore semantics quickly when the architecture calls for it.

On AI tooling specifically, I deployed a \textbf{Retrieval-Augmented Generation (RAG)} pipeline with an \textbf{agentic workflow} at Amazon that autonomously retrieves operational logs to reason over incident details, and implemented a \textbf{Model Context Protocol (MCP) server} for tool-augmented retrieval. Bringing ML and agent patterns into production infrastructure isn't a side experiment for me, it's been the shape of my last year of work. The fraud-detection-with-Kafka-Streams direction the role describes is exactly where I'd want to ramp.

What pulls me toward Telnyx is the seriousness of the infrastructure. Telecommunications and messaging at global scale runs on real consequences: every dropped message is a failure for someone, every region has to behave the same way, and uptime is measured in nines that matter. That's the kind of feedback loop I want to be inside of, alongside engineers pushing code to production multiple times a day.

I'd welcome the chance to discuss further. Thanks for considering my application.

\vspace{12pt}

Sincerely, \\
Ajay Venkatesh

\end{document}
